% GENERATED FILE. Edit canonical lessons, then run npm run generate:books.
% canonical-source-hash: fnv1a64:6464dde02ae01fd1
% canonical-lessons: RU-C208-chuvstvovat, RU-C208-zamechat, RU-C208-slyshat, RU-C208-nyukhat, RU-C208-trogat

\chapter{To feel, To notice, To hear, To sniff, To touch}
\label{ch:ru-to-feel-to-notice-to-hear-to-sniff-to-touch}
\hlchaptermodality{208}

\begin{chapteropening}
\textbf{By the end of this chapter:} \emph{I can say to feel, to notice, to hear, to sniff and to touch.}

Complete the last lesson of chapter 208: I can say to feel, to notice, to hear, to sniff and to touch.
\end{chapteropening}

\section[\texorpdfstring{\ru{чувствовать} (chúvstvovat)}{чувствовать (chúvstvovat)}]{\ru{чувствовать} (chúvstvovat) — to feel}

\label{lesson:RU-C208-chuvstvovat}



\begin{quote}
Before the new one: say the Russian for to cough, then the Russian for to sneeze.
\end{quote}

\subsection*{You'll want to know: \ru{чувствовать}}
\textbf{\ru{чувствовать}} (\emph{chúvstvovat}) — \textquotedblleft{}to feel\textquotedblright{}.

\begin{grammarlens}[title={What you've built}]
One more word for this chapter.
\end{grammarlens}

\subsection*{Your turn}
\begin{itemize}
\raggedright
  \item \emph{Say it:} \emph{chúvstvovat}
  \item \emph{Say it:} \emph{chúvstvovat}, once more
  \item \emph{Say it:} say \emph{chúvstvovat}
  \item \emph{Recall:} say the Russian for to cough, then the Russian for to sneeze, then say \emph{chúvstvovat} again
\end{itemize}

\begin{tcolorbox}[breakable,colback=teal!4,colframe=teal!35!black,title={Before you move on}]
What does \ru{чувствовать} mean? (\textquotedblleft{}To feel\textquotedblright{}.) Say it once more.
\end{tcolorbox}
\section[\texorpdfstring{\ru{замечать} (zamechát)}{замечать (zamechát)}]{\ru{замечать} (zamechát) — to notice}

\label{lesson:RU-C208-zamechat}



\begin{quote}
Before the new one: say the Russian for to sneeze, then the Russian for to feel.
\end{quote}

\subsection*{You'll want to know: \ru{замечать}}
\textbf{\ru{замечать}} (\emph{zamechát}) — \textquotedblleft{}to notice\textquotedblright{}.

\begin{grammarlens}[title={What you've built}]
One more word for this chapter.
\end{grammarlens}

\subsection*{Your turn}
\begin{itemize}
\raggedright
  \item \emph{Say it:} \emph{zamechát}
  \item \emph{Say it:} \emph{zamechát}, once more
  \item \emph{Say it:} say \emph{zamechát}
  \item \emph{Recall:} say the Russian for to sneeze, then the Russian for to feel, then say \emph{zamechát} again
\end{itemize}

\begin{tcolorbox}[breakable,colback=teal!4,colframe=teal!35!black,title={Before you move on}]
What does \ru{замечать} mean? (\textquotedblleft{}To notice\textquotedblright{}.) Say it once more.
\end{tcolorbox}
\section[\texorpdfstring{\ru{слышать} (slýshat)}{слышать (slýshat)}]{\ru{слышать} (slýshat) — to hear}

\label{lesson:RU-C208-slyshat}



\begin{quote}
Before the new one: say the Russian for to feel, then the Russian for to notice.
\end{quote}

\subsection*{You'll want to know: \ru{слышать}}
\textbf{\ru{слышать}} (\emph{slýshat}) — \textquotedblleft{}to hear\textquotedblright{}.

\begin{grammarlens}[title={What you've built}]
One more word for this chapter.
\end{grammarlens}

\subsection*{Your turn}
\begin{itemize}
\raggedright
  \item \emph{Say it:} \emph{slýshat}
  \item \emph{Say it:} \emph{slýshat}, once more
  \item \emph{Say it:} say \emph{slýshat}
  \item \emph{Recall:} say the Russian for to feel, then the Russian for to notice, then say \emph{slýshat} again
\end{itemize}

\begin{tcolorbox}[breakable,colback=teal!4,colframe=teal!35!black,title={Before you move on}]
What does \ru{слышать} mean? (\textquotedblleft{}To hear\textquotedblright{}.) Say it once more.
\end{tcolorbox}
\section[\texorpdfstring{\ru{нюхать} (nyúkhat)}{нюхать (nyúkhat)}]{\ru{нюхать} (nyúkhat) — to sniff, to smell}

\label{lesson:RU-C208-nyukhat}



\begin{quote}
Before the new one: say the Russian for to notice, then the Russian for to hear.
\end{quote}

\subsection*{You'll want to know: \ru{нюхать}}
\textbf{\ru{нюхать}} (\emph{nyúkhat}) — \textquotedblleft{}to sniff, to smell\textquotedblright{}.

\begin{grammarlens}[title={What you've built}]
One more word for this chapter.
\end{grammarlens}

\subsection*{Your turn}
\begin{itemize}
\raggedright
  \item \emph{Say it:} \emph{nyúkhat}
  \item \emph{Say it:} \emph{nyúkhat}, once more
  \item \emph{Say it:} say \emph{nyúkhat}
  \item \emph{Recall:} say the Russian for to notice, then the Russian for to hear, then say \emph{nyúkhat} again
\end{itemize}

\begin{tcolorbox}[breakable,colback=teal!4,colframe=teal!35!black,title={Before you move on}]
What does \ru{нюхать} mean? (\textquotedblleft{}To sniff, to smell\textquotedblright{}.) Say it once more.
\end{tcolorbox}
\section[\texorpdfstring{\ru{трогать} (trógat)}{трогать (trógat)}]{\ru{трогать} (trógat) — to touch}

\label{lesson:RU-C208-trogat}



\begin{quote}
Before the new one: say the Russian for to hear, then the Russian for to sniff.
\end{quote}

\subsection*{You'll want to know: \ru{трогать}}
\textbf{\ru{трогать}} (\emph{trógat}) — \textquotedblleft{}to touch\textquotedblright{}.

\begin{grammarlens}[title={What you've built}]
One more word for this chapter.
\end{grammarlens}

\subsection*{Your turn}
\begin{itemize}
\raggedright
  \item \emph{Say it:} \emph{trógat}
  \item \emph{Say it:} \emph{trógat}, once more
  \item \emph{Say it:} say \emph{trógat}
  \item \emph{Recall:} say the Russian for to hear, then the Russian for to sniff, then say \emph{trógat} again
\end{itemize}

\begin{tcolorbox}[breakable,colback=teal!4,colframe=teal!35!black,title={Before you move on}]
What does \ru{трогать} mean? (\textquotedblleft{}To touch\textquotedblright{}.) Say it once more.
\end{tcolorbox}
